% ch3_b 第3章 §3.2尾–§3.3 (书页 79–89 / p094–p104)
% ——— 块起点接续说明 ———
% 书页 78（ch3_a）末行止于 "The energy density (or, more precisely, the"，
% 本块自书页 79 顶部的半句续起（"density of the thermal potential) of the
% classical ground state is"）：
%JOIN%
\begin{equation*}
\frac{\Omega_0}{\mathcal{V}} = -\mu|\varphi_0|^2 + \frac{V_0}{2}|\varphi_0|^4
\end{equation*}

%FIG{3.3}: {序参量 $\varphi_0$ 与超流相变。}

对能量求极小，我们看到，对于 $\mu < 0$，基态由 $\varphi_0 = 0$ 表征，对应于没有玻色子的态。对于 $\mu > 0$，基态是简并的，由玻色场的下列幅值表征：
\begin{equation*}
\varphi_0 = \sqrt{\frac{\mu}{V_0}}\,\mathrm{e}^{\mathrm{i}\theta}\tag{3.3.3}
\end{equation*}
其中 $\theta$ 可任取。这样的态具有玻色子密度 $\rho_0 = \mu/V_0$。热势作为 $\mu$ 的函数由下式给出：
\begin{equation*}
\Omega_0(\mu) =
\begin{cases}
0, & \text{当 } \mu < 0\\[6pt]
\dfrac{\mu^2}{V_0}, & \text{当 } \mu > 0
\end{cases}
\end{equation*}
（译注：由上一式可得 $\mu>0$ 时 $\Omega_0/\mathcal{V} = -\mu^2/(2V_0)$；原书此处印作 $\mu^2/V_0$，未含体积 $\mathcal{V}$ 与因子 $-1/2$，疑为排印之误。）

我们看到，作为 $\mu$ 的函数，基态能量 $\Omega_0(\mu)$ 在 $\mu = 0$ 处有一个奇点。该奇点标志着 $\mu = 0$ 处的一个二级相变，因为 $\Omega_0''(\mu)$ 在该点有一跳变（见图 3.2）。这个相变称为超流相变（superfluid transition）。\footnote{我们将在 3.7.3 节讨论超流性。}

超流相的特征是非零的玻色场 $\varphi(\pmb{x}) \neq 0$，而无玻色子相的特征是玻色场为零 $\varphi(\pmb{x}) = 0$。超流相变由 $\varphi_0 = 0$ 到 $\varphi_0 \neq 0$ 的改变所标志（见图 3.3）。因此，我们称 $\varphi_0$ 为超流的序参量。

超流相还以对称性破缺为特征。我们看到，我们的哈密顿量或拉格朗日量具有 $U(1)$ 对称性，因为在变换
\begin{equation*}
\varphi \rightarrow \mathrm{e}^{\mathrm{i}\theta}\varphi\tag{3.3.4}
\end{equation*}
之下，无论我们处于哪个相，它们都保持不变。然而，超流相中的（经典）基态（见式 (3.3.3)）不尊重 $U(1)$ 对称性，也就是说，它们在 $U(1)$ 变换 (3.3.4) 下不是不变的。

%FIG{3.4}: {（a）与（b）：在不同极小值之间的切换导致（c）中 $\Omega_0$ 的一个拐折，它代表一个一级相变。}

%FIG{3.5}: {在存在 $\phi_0 \to -\phi_0$ 对称性时，极小值之间的切换可以连续地进行，并导致一个二级相变。（a）基态在 $\phi_0 \to -\phi_0$ 条件下是对称的。（b）转变点。（c）基态破坏 $\phi_0 \to -\phi_0$ 对称性。}

超流相变展示了一条深刻的原理：以热势 $\Omega$ 的非解析点为特征的连续相变，与态的对称性的改变相关联。为领会这一点，让我们从头说起。我们知道，相变是由系统的能量或自由能中的奇点引起的。所以，要理解相变，就需要理解（自由）能是如何产生奇点的。

对于相互作用玻色系统，能量函数 $\Omega_0(\mu; \phi_0)$ 作为 $\phi_0$ 和 $\mu$ 的函数是一个没有奇点的解析函数。基态能量 $\Omega_0(\mu)$ 是 $\Omega_0(\mu; \phi_0)$ 对 $\phi_0$ 取的极小值。一般地，$\Omega_0(\mu)$ 也是解析函数。然而，有时 $\Omega_0(\mu; \phi_0)$ 作为 $\phi_0$ 的函数有多个局部极小。当我们改变 $\mu$ 时，全局极小可能从一个局部极小切换到另一个（见图 3.4）。在这种情况下，所得的 $\Omega_0(\mu)$ 将在切换点有一个奇点。这样的奇点代表一个一级相变，因为 $\Omega_0'(\mu)$ 有一跳变（见图 3.4(c)）。

然而，当能量函数 $\Omega_0(\mu,\phi_0)$ 具有对称性，比如说 $\Omega_0(\mu,\phi_0) = \Omega_0(\mu,\mathrm{e}^{\mathrm{i}\theta}\phi_0)$ 时，极小值之间的切换可以有非常不同的行为，并且可以是连续的（见图 3.5）。由于这一对称性，新的全局极小是简并的。系统面临一个艰难的抉择：要挑哪一个极小驻留。系统做出选择之后，新相便不再具有 $\phi_0 \to \mathrm{e}^{\mathrm{i}\theta}\phi_0$ 对称性。这种从对称系统中演生出非对称态的现象称为自发对称性破缺。我们看到，二级相变与基态对称性的改变密切相关。这正是 Landau 相与相变对称性破缺理论的核心（Landau, 1937）。这一简单的思想应用极为广泛。实际上，几乎所有连续相变都以某种自发对称性破缺为特征。

一般地，自发对称性破缺可以用一个序参量来表征。序参量可以选为一个在对称变换下非平凡变换的算符的期望值。在上面，我们选 $\langle\varphi\rangle$ 作为序参量。由于其变换性质 $\langle\varphi\rangle \to \mathrm{e}^{\mathrm{i}\theta}\langle\varphi\rangle$，在对称相中 $\langle\varphi\rangle = 0$。相比之下，在对称性破缺相中 $\langle\varphi\rangle = \varphi_0 \neq 0$。

在对称性破缺相中，如果我们知道 $\varphi$ 场在一点的相位，那么我们就知道它在任何其他点的相位。这一性质称为长程序。$\varphi$ 场中的长程关联可以在数学上表示为
\begin{equation*}
\left.\left\langle \varphi(t,\pmb{x})^{\dagger}\varphi(0,0)\right\rangle\right|_{(t,\pmb{x})\to\infty} = |\varphi_0|^2 \neq 0
\end{equation*}
如果 $\varphi$ 的相位在各点之间是无规且互不关联的，那么上述关联将为零。当然，这里我们只考虑没有量子涨落的经典基态。在这种情形下我们看到，非零序参量与 $\varphi$ 算符中的长程序密切相关。对称性破缺相既可以用非零序参量表征，也可以用长程关联表征。在下面几节中，我们将把量子涨落包括进来，看看它们会如何影响序参量和长程序。

\subsection*{问题 3.3.1}

由作用量 (3.3.1) 导出经典运动方程。只考虑基态附近的小涨落，对对称相（$\mu < 0$）与对称性破缺相（$\mu > 0$）分别导出线性化的运动方程。证明对称基态附近的涨落具有有限能隙，而对称性破缺基态附近的涨落包含一个零能隙的模式。此模式具有线性色散 $\omega \propto k$。

\subsection*{问题 3.3.2}

水–水蒸气相变可以用下列吉布斯能（Gibbs energy）函数描述：
\begin{equation*}
G(T,P;n) = h(T,P)n + a(T,P)n^2 + c(T,P)n^3 + b(T,P)n^4
\end{equation*}
其中 $n$ 是水分子的密度。系统的吉布斯能 $G(T,P)$ 由 $G(T,P;n)$ 对 $n$ 取极小值给出。系统有一条一级相变线，它终止于一个临界点。试求确定一级相变线与临界点的方程。（提示：先考虑 $c = 0$ 的情形。如果把 $n$ 看作自旋矩的密度，这样的吉布斯能与磁场 $\delta\rho$ 中伊辛模型的自由能具有相同的形式。）

\subsection{低能有效理论}

\begin{keypoints}
\item 描述慢（即低能）涨落的有效拉格朗日量，可以通过积掉快（即高能）涨落得到。
\end{keypoints}

为研究经典基态之上的激发，让我们考虑经典基态附近的下述小涨落：
\begin{equation*}
\varphi = \varphi_0 + \delta\varphi
\end{equation*}
$\delta\varphi$ 的动力学由同一个作用量 (3.3.1) 支配。但现在我们可以假设 $\delta\varphi$ 很小，只保留 $\delta\varphi$ 的二次项。对于对称相 $\varphi_0 = 0$，我们有
\begin{equation*}
S = \int \mathrm{d}^d\pmb{x}\,\mathrm{d}t\,\Big[\mathrm{i}\frac{1}{2}(\varphi^*\partial_t\varphi - \varphi\partial_t\varphi^*) - \frac{1}{2m}\partial_{\pmb{x}}\varphi^*\,\partial_{\pmb{x}}\varphi + \mu\varphi^*\varphi\Big]\tag{3.3.5}
\end{equation*}
在对称相中，运动方程为
\begin{equation*}
\Big[\mathrm{i}\partial_t - \frac{1}{2m}(-\mathrm{i}\partial_{\pmb{x}})^2 + \mu\Big]\varphi = 0
\end{equation*}
由此得到涨落的如下色散：
\begin{equation*}
\omega = \frac{1}{2m}k^2 - \mu\tag{3.3.6}
\end{equation*}
在下一节中我们将看到，这些涨落对应于类粒子激发：涨落的频率与波矢量分别变成粒子的能量与动量。我们看到，激发具有有限能隙 $\Delta = -\mu$（注意 $\mu < 0$），这正是产生一个激发所需的能量。事实上，由色散 (3.3.6) 所描述的粒子恰好对应于构成我们系统的那些玻色子。

为理解对称性破缺相中的低能激发，我们希望推导一个只包含与低能激发相关的模式的低能有效理论。至于哪些涨落具有低能量，我们注意到，在对称性破缺相中，凝聚体的不同相位 $\theta$ 导致简并的基态。如果我们在一个很大但有限的区域内改变凝聚体的相位，那么系统在局部上仍处于它的某个简并基态。系统得以知道自己处于激发态的唯一途径是通过 $\theta$ 的梯度。如果相位在空间中变化缓慢，相应的激发就具有较小的能量。因此，低能激发对应于诸简并基态之间的涨落，由 $\theta$ 场描述。基于这一图像，我们写出
\begin{equation*}
\varphi = \sqrt{\rho_0 + \delta\rho}\,\mathrm{e}^{\mathrm{i}\theta}\tag{3.3.7}
\end{equation*}
其中 $\rho_0 = \varphi_0^2 = \mu/V_0$ 是基态的密度，$\delta\rho$ 是密度涨落。在这种写法中，$\theta$ 描述低能的慢涨落，$\delta\rho$ 描述高能的快涨落。把式 (3.3.7) 代入式 (3.3.1)，并展开到 $\theta$ 与 $\delta\rho$ 的二次阶，我们得到
\begin{equation*}
S = \int \mathrm{d}^d\pmb{x}\,\mathrm{d}t\,\Big[-(\rho_0+\delta\rho)\partial_t\theta - \frac{\rho_0(\partial_{\pmb{x}}\theta)^2}{2m} - \frac{(\partial_{\pmb{x}}\delta\rho)^2}{8m\rho_0} - \frac{V_0}{2}\delta\rho^2\Big]\tag{3.3.8}
\end{equation*}

这里我们面对的是多体物理中一个十分典型的问题。我们有两个场，一个描述慢的低频（即低能）涨落，另一个描述快的高频涨落。如果我们只关心低能物理，如何移除快场，以得到一个简单的低能有效理论？

得到低能有效理论的一种办法是令 $\delta\rho = 0$，把 $\delta\rho$ 场丢掉。这导致 $\theta$ 的下列低能有效作用量：
\begin{equation*}
S = \int \mathrm{d}^d\pmb{x}\,\mathrm{d}t\,\Big[-\varphi_0^2\partial_t\theta - \frac{1}{2m}\varphi_0^2(\partial_{\pmb{x}}\theta)^2\Big]
\end{equation*}
然而，这个有效作用量并不十分正确，因为 $\partial_t\theta$ 是一个全导数（total derivative），它不进入运动方程。我们需要得到 $(\partial_t\theta)^2$ 项，才能理解 $\theta$ 的低能动力学性质。要得到 $(\partial_t\theta)^2$ 项，移除 $\delta\rho$ 场时就必须更加小心。

得到低能有效理论的一个更好的办法，是在路径积分中积分掉快场 $\delta\rho$。在这里这很容易做到，因为作用量 (3.3.8) 是 $\delta\rho$ 的二次型。完成高斯积分后，我们发现
\begin{align*}
Z &= \int \mathcal{D}\theta\,\mathrm{e}^{\mathrm{i}\int \mathrm{d}^d\pmb{x}\,\mathrm{d}t\,\big[-\frac{\rho_0}{2m}(\partial_{\pmb{x}}\theta)^2 + \frac{1}{2}\partial_t\theta\,\frac{1}{V_0 - \frac{1}{4m\rho_0}\partial_{\pmb{x}}^2}\,\partial_t\theta\big]}\\
&\approx \int \mathcal{D}\theta\,\mathrm{e}^{\mathrm{i}\int \mathrm{d}^d\pmb{x}\,\mathrm{d}t\,\big[-\frac{\rho_0}{2m}(\partial_{\pmb{x}}\theta)^2 + \frac{1}{2V_0}(\partial_t\theta)^2\big]}\tag{3.3.9}
\end{align*}
其中我们假设了 $\theta$ 场在空间中变化缓慢，并丢掉了 $\frac{1}{4m\rho_0}(\partial_{\pmb{x}})^2$ 项；我们还丢掉了作用量中的全时间导数项。\footnote{在存在涡旋时，全时间导数项是重要的，那时不应将其丢掉。见 3.6 节。}我们得到 $\theta$ 的下列简单的低能有效作用量：
\begin{equation*}
S_{\mathrm{eff}} = \int \mathrm{d}^d\pmb{x}\,\mathrm{d}t\,\Big[\frac{1}{2V_0}(\partial_t\theta)^2 - \frac{\rho_0}{2m}(\partial_{\pmb{x}}\theta)^2\Big]\tag{3.3.10}
\end{equation*}
注意 $\theta$ 场实际上生活在一个圆上，因此 $\theta$ 与 $\theta + 2\pi$ 描述同一点。为更准确起见，我们可以引入 $z = \mathrm{e}^{\mathrm{i}\theta}$，把上式改写为
\begin{equation*}
S_{\mathrm{eff}} = \int \mathrm{d}^d\pmb{x}\,\mathrm{d}t\,\Big[\frac{1}{2V_0}|\partial_t z|^2 - \frac{\rho_0}{2m}|\partial_{\pmb{x}} z|^2\Big]\tag{3.3.11}
\end{equation*}
由式 (3.3.10) 或式 (3.3.11) 所描述的模型称为 XY 模型（XY-model）。

经典基态由均匀的 $\theta$ 场给出：$\theta(\pmb{x},t) = $ 常数。激发由 $\theta$ 的涨落描述，这些涨落满足运动方程
\begin{equation*}
\Big(-\frac{1}{2V_0}\partial_t^2 + \frac{\rho_0}{2m}\partial_{\pmb{x}}^2\Big)\theta = 0
\end{equation*}
上式是一个波动方程，它描述一列具有线性色散的波
\begin{equation*}
\omega = v|\pmb{k}|,\qquad v^2 = \frac{\rho_0 V_0}{m}\tag{3.3.12}
\end{equation*}
其中 $v$ 是波速。

在许多情形下，仅仅知道低能有效作用量是不够的。同样重要的是要知道，原理论中的场（或算符）如何由低能有效理论中的场（或算符）来表示。下面我们以密度算符为例，说明原理论中的物理算符（或场）如何由有效理论中的场来表示。首先，在原拉格朗日量中加一个与密度耦合的源项 $-A_0\rho = -A_0(\rho_0 + \delta\rho)$。然后，进行同样的计算以得到有效理论。我们发现，有效理论中含有一个附加项（准确到 $A_0$ 的线性阶）$-A_0(\rho_0 - \frac{\partial_t\theta}{V_0})$。于是，密度算符在有效理论中表示为
\begin{equation*}
\rho = \rho_0 - \frac{\partial_t\theta}{V_0}\tag{3.3.13}
\end{equation*}
类似地，我们发现玻色子流密度 $\pmb{j} = \mathrm{Re}\,\varphi^{\dagger}\frac{\partial_{\pmb{x}}}{\mathrm{i}m}\varphi$ 在有效理论中变为
\begin{equation*}
\pmb{j} = \frac{\rho_0}{m}\partial_{\pmb{x}}\theta\tag{3.3.14}
\end{equation*}
式 (3.3.13)、式 (3.3.14) 与式 (3.3.10) 使我们能够在简单的低能有效理论之内，用路径积分计算密度关联与流关联。这些关联随后便可与诸如压缩率（compressibility）和电导率等实验结果相比较。$\rho$ 与 $\theta$ 之间的关系还告诉我们，$\theta$ 所描述的低能涨落正是密度–流涨落。

\subsection*{问题 3.3.3}

\begin{enumerate}
\item[(a)] 得到 $\theta$ 的低能有效理论还有另一种办法：通过求解 $\delta\rho$ 场的运动方程，把 $\delta\rho$ 用 $\partial_t\theta$、$\partial_{\pmb{x}}\theta$ 等表示出来；然后把 $\delta\rho$ 代回作用量，得到一个只含 $\theta$ 的有效作用量。证明这一方法给出同样的 XY 模型作用量 (3.3.10)。
\item[(b)] 如果把 $\delta\rho$ 代入密度算符和流算符，则会分别重新得到式 (3.3.13) 与式 (3.3.14)。证明 $\rho$ 和 $\pmb{j}$ 满足守恒定律 $\partial_t\rho + \partial_{\pmb{x}}\cdot\pmb{j} = 0$。
\item[(c)] 用式 (3.3.13) 给出的密度算符表达式，计算动量–频率空间中超流的密度–密度关联函数。
\end{enumerate}

\subsection{波即是粒子}

\begin{keypoints}
\item 超流态中演生出一种新型的粒子——准粒子（quasiparticle）。准粒子是玻色型的、无相互作用的，并且具有线性色散。
\end{keypoints}

我们一直把经典基态附近的 $\theta$ 涨落（或者对称相中的 $\delta\varphi$）看作波。由于量子物理中的波粒二象性（particle–wave duality），波也可以看作粒子。这些粒子称为准粒子。下面我们将量子化低能有效理论，得到准粒子的量子理论。该量子理论结果是一个玻色子理论，这告诉我们准粒子是玻色子。

我们先把低能有效拉格朗日量 (3.3.10) 在 $\pmb{k}$ 空间中写成如下形式：
\begin{equation*}
L_{\mathrm{eff}} = \sum_{\pmb{k}}\Big[\frac{A_{\pmb{k}}}{2}\dot{\theta}_{-\pmb{k}}\dot{\theta}_{\pmb{k}} - \frac{B_{\pmb{k}}}{2}\theta_{-\pmb{k}}\theta_{\pmb{k}}\Big],\tag{3.3.15}
\end{equation*}
其中
\begin{equation*}
A_{\pmb{k}} = V_0^{-1},\qquad B_{\pmb{k}} = \frac{\rho_0 k^2}{m},\tag{3.3.16}
\end{equation*}
$\theta_{\pmb{k}} = \mathcal{V}^{-1/2}\int \mathrm{d}^d\pmb{x}\,\theta(\pmb{x})\mathrm{e}^{-\mathrm{i}\pmb{k}\cdot\pmb{x}}$，而 $\mathcal{V}$ 是系统的总体积。上述作用量描述一组彼此退耦的谐振子。

为得到描述量子化理论的量子哈密顿量，我们先算出经典哈密顿量。$\theta_{\pmb{k}}$ 的正则动量（canonical momentum）为
\begin{equation*}
\pi_{\pmb{k}} = \frac{\partial L_{\mathrm{eff}}}{\partial \dot{\theta}_{\pmb{k}}} = A_{\pmb{k}}\dot{\theta}_{-\pmb{k}}
\end{equation*}
经典哈密顿量 $H = \sum_{\pmb{k}}\pi_{\pmb{k}}\dot{\theta}_{\pmb{k}} - L_{\mathrm{eff}}$ 由下式给出：
\begin{equation*}
H = \sum_{\pmb{k}}\Big[\frac{1}{2A_{\pmb{k}}}\pi_{-\pmb{k}}\pi_{\pmb{k}} + \frac{B_{\pmb{k}}}{2}\theta_{-\pmb{k}}\theta_{\pmb{k}}\Big]
\end{equation*}
只需把 $\pi_{\pmb{k}}$ 和 $\theta_{\pmb{k}}$ 当作满足下列正则对易关系（canonical commutation relation）的算符，就得到量子哈密顿量：
\begin{equation*}
[\theta_{\pmb{k}}, \pi_{\pmb{k}'}] = \mathrm{i}\delta_{\pmb{k}\pmb{k}'}
\end{equation*}
现在我们引入
\begin{equation*}
\alpha_{\pmb{k}} = u_{\pmb{k}}\theta_{\pmb{k}} + \mathrm{i}v_{\pmb{k}}\pi_{-\pmb{k}},\quad u_{\pmb{k}} = \frac{1}{\sqrt{2}}(A_{\pmb{k}}B_{\pmb{k}})^{1/4},\quad v_{\pmb{k}} = \frac{1}{\sqrt{2}}(A_{\pmb{k}}B_{\pmb{k}})^{-1/4}.\tag{3.3.17}
\end{equation*}
我们发现 $\alpha_{\pmb{k}}$ 满足振子下降算符（lowering operator）的代数：
\begin{equation*}
[\alpha_{\pmb{k}}, \alpha_{\pmb{k}'}^{\dagger}] = \delta_{\pmb{k}\pmb{k}'},\qquad [\alpha_{\pmb{k}}, \alpha_{\pmb{k}'}] = 0.
\end{equation*}
用下降与上升算符（raising operator）$(\alpha_{\pmb{k}}, \alpha_{\pmb{k}}^{\dagger})$ 表示，哈密顿量取振子的标准形式：
\begin{equation*}
H = \sum_{\pmb{k}}\epsilon_{\pmb{k}}\alpha_{\pmb{k}}^{\dagger}\alpha_{\pmb{k}} + \text{常数},\qquad \epsilon_{\pmb{k}} = \sqrt{\frac{B_{\pmb{k}}}{A_{\pmb{k}}}} = \sqrt{\frac{V_0\rho_0}{m}}\,|\pmb{k}| = v|\pmb{k}|\tag{3.3.18}
\end{equation*}
这正是具有单玻色子能量 $\epsilon_{\pmb{k}} = v|\pmb{k}|$ 的自由玻色子哈密顿量 (3.1.1)。

我们看到，$\theta$ 的集体涨落产生了具有线性色散的玻色型准粒子。由于 $\theta$ 的涨落对应于超流体中的密度波，我们也可以说，密度波在量子理论中变成了离散的准粒子。这一图像适用于更普遍的情形。事实上，量子基态的任何波状涨落在量子化之后都对应于离散的准粒子。\footnote{空气中的声波不对应于任何离散的准粒子。这是因为声波不是任何量子基态的涨落，因此它不对应于基态之上的任何激发。}固体中的声波变成由哈密顿量 (3.3.18) 描述的声子（phonon）。因此，我们也将把超流体中的准粒子称作声子。声子速度（即密度波的速度）$v = \sqrt{V_0\rho_0/m}$ 与平均场结果 (3.2.9) 相一致。

我们想强调，玻色型准粒子——声子——与构成超流体的原来那些玻色子截然不同。原来的玻色子具有二次色散 $\epsilon_{\pmb{k}} = k^2/2m$，而声子具有线性色散 $\epsilon_{\pmb{k}} = v|\pmb{k}|$；原来的玻色子是有相互作用的，而声子是自由的。事实上，一个声子对应于许多原来玻色子的一种集体运动。这是从相互作用玻色系统中演生出一种新型玻色子——具有线性色散的无相互作用声子——的一个例子。

\subsection{作为玩具宇宙的超流体}

\begin{keypoints}
\item 如果我们把超流体看作一个玩具宇宙，那么声子将是这个玩具宇宙中无质量的“基本粒子”。
\item 旋子（roton，即涡环 vortex ring）可以通过交换声子发生相互作用。这导致两个旋子之间的 $1/r^4$ 偶极相互作用。
\end{keypoints}

为了领会我们的宇宙以及其中的基本粒子，让我们想象一个由超流体形成的玩具宇宙（toy universe）。假设构成超流体的玻色子太小，玩具宇宙中的居民看不见它们。问题是：在玩具宇宙的居民看来，这个玩具宇宙是什么样子的？

%FIG{3.6}: {（a）旋子（涡环）周围与（b）声子电荷的偶极子周围的玻色子流动图案。}

这个玩具宇宙包含一种无质量激发\footnote{质量为 $m$ 的相对论性粒子具有色散 $\epsilon_{\pmb{k}} = \sqrt{m^2c^4 + c^2\pmb{k}^2}$，其中 $c$ 是光速。声子色散 $\epsilon_{\pmb{k}} = v|\pmb{k}|$ 可以看作 $\sqrt{m^2c^4 + c^2\pmb{k}^2}$ 的无质量极限，其中声子速度 $v$ 扮演光速的角色。}——声子。对玩具宇宙中的居民来说，声子只是一种类粒子激发，因为没有人能看见构成超流体的玻色子，也没有人知道声子从何而来。于是，玩具宇宙的居民便把声子称作一种“基本粒子”。

我们看到，这个玩具宇宙与我们的宇宙十分相似。我们的宇宙也有一种无质量激发——光子。我们也把光子称作“基本粒子”。与声子一样，光子彼此之间也不相互作用。然而，电荷之间的 $1/r^2$ 库仑定律正是由光子传递的。那么，声子是否也会在它们的“电荷”之间诱导出类似的相互作用呢？

上述问题的答案是肯定的。不过，我们首先需要解释声子的“电荷”是什么。我们知道，电荷产生守恒的通量——电场。类似地，声子的“电荷”也产生守恒的通量，只不过对声子而言，这个通量是超流体中玻色子的通量。按定义，一个正的“声子电荷”（phonon charge）是玻色子的源：玻色子在某个点以恒定速率被产生；一个负的声子电荷则是玻色子的汇：玻色子在某个点以恒定速率被湮没。

为理解声子电荷之间的相互作用，我们对满足约束 $\partial_{\pmb{x}}\cdot\pmb{j} \allowbreak- J^0 = 0$ 的 $\theta$ 极小化能量 $\int \mathrm{d}^3\pmb{x}\,\allowbreak\frac{\rho_0}{2m}(\partial_{\pmb{x}}\theta)^2$。这里 $\pmb{j} = \frac{\rho_0}{m}\partial_{\pmb{x}}\theta$ 是玻色子流，$J^0(\pmb{x},t) = \sum_i q_i\delta(\pmb{x} - \pmb{x}_i)$ 是声子电荷的密度，$q_i$ 是第 $i$ 个声子电荷的值。$q_i > 0$ 对应于源，$q_i < 0$ 对应于汇。我们发现，两个声子电荷 $q_1$ 与 $q_2$ 之间的势正比于 $q_1q_2/r$。声子电荷之间的相互作用与库仑定律非常相似：力正比于 $1/r^2$，同号电荷相斥，异号电荷相吸。

在真实的超流体中，玻色子是守恒的。因此，源和汇——亦即声子电荷——是不允许存在的。唯一允许的激发是低能声子与高能旋子（见图 3.1）。应当注意，旋子可以看作由涡线（vortex line）构成的一个小环。\footnote{$(2+1)$ 维超流体中的涡旋定义见 3.4.2 节。这一定义可以推广到任意维度。}如果做多极展开（multipole expansion），旋子周围玻色子的流动图案可以看作由一组声子电荷产生的流动图案。玻色子数守恒要求多极展开中的声子电荷总和为零。涡环的对称性允许多极展开中存在有限的偶极矩（见图 3.6），而旋子的偶极矩一般不为零。因此，旋子彼此之间存在 $1/r^4$ 偶极相互作用。

总而言之，这个玩具宇宙包含带有偶极矩的粒子，其长程偶极相互作用通过交换无质量的声子而产生。

\subsection{自发对称性破缺与无能隙激发}

\begin{keypoints}
\item 普适性质（universal properties）的概念。
\item 连续对称性的自发破缺总是导致无能隙的玻色型激发。
\end{keypoints}

我们想强调，在超流体中，声子的三个性质——无相互作用、无能隙的线性色散以及玻色统计——是普适性质，不依赖于玻色子之间相互作用的细节。在上面，我们假设玻色子具有 $\delta$ 函数相互作用。然而，只要相互作用 $V(\pmb{x}-\pmb{x}')$ 是短程的，即使玻色子具有更一般的相互作用 $L_{int} = -\int \mathrm{d}\pmb{x}\,\mathrm{d}\pmb{x}'\,\frac{1}{2}|\varphi(\pmb{x})|^2 V(\pmb{x}-\pmb{x}')|\varphi(\pmb{x}')|^2$，这些普适性质也保持不变。既然准粒子的无能隙性是不依赖于初始拉格朗日量细节的普适性质，就应当存在着一种不依赖这些细节的关于无能隙性的一般性理解。事实表明，超流相中无能隙激发的存在，与拉格朗日量的 $U(1)$ 对称性以及基态的自发对称性破缺密切相关。下面我们将解释这一关系。

拉格朗日量 (3.3.1) 具有 $U(1)$ 对称性，因为它在 $U(1)$ 变换 $\varphi \to \mathrm{e}^{\mathrm{i}\eta}\varphi$ 下不变。如果我们显式破缺这一对称性，即使拉格朗日量不再尊重 $U(1)$ 对称性，低能模式也将获得有限能隙。例如，若在拉格朗日量中加入 $C\,\mathrm{Re}\varphi$ 项以显式破缺 $U(1)$ 对称性，则 XY 模型中会诱导出 $\cos(\theta)$ 形式的项：$\mathcal{L} = \frac{1}{2V_0}(\partial_t\theta)^2 - \frac{\rho_0}{2m}(\partial_{\pmb{x}}\theta)^2 + g\cos(\theta)$。经典基态由 $\theta = 0$ 给出。对于基态附近的小涨落，拉格朗日量可以简化为 $\mathcal{L} = \frac{1}{2V_0}(\partial_t\theta)^2 - \frac{\rho_0}{2m}(\partial_{\pmb{x}}\theta)^2 - \frac{g}{2}\theta^2$。所得的波动方程 $(-\frac{1}{2V_0}\partial_t^2 + \frac{\rho_0}{2m}\partial_{\pmb{x}}^2 - g)\theta = 0$ 告诉我们，无论波矢量如何，波总有有限频率 $\omega = \sqrt{2V_0g + \frac{V_0\rho_0}{m}k^2}$。于是，势项 $\cos(\theta)$ 为 $\theta$ 准粒子打开了一个能隙。

如果拉格朗日量和基态在 $U(1)$ 变换下都不变，那么对称性就没有破缺。由式 (3.3.6) 我们看到，对称基态一般也具有有限能隙。

%FIG{3.7}: {无能隙的 Nambu–Goldstone 模是简并基态之间的涨落。}

超流相自发破缺 $U(1)$ 对称性：拉格朗日量具有 $U(1)$ 对称性，而基态没有（见 3.3.2 节）。只有在这种情形下，无能隙的声子激发才存在。

Nambu 和 Goldstone 证明了一条普遍定理：如果一个连续对称性在一个相中被\emph{自发}破缺，那么该相必定包含无能隙激发（Nambu, 1960; Goldstone, 1961）。这些无能隙激发通常称为 Nambu–Goldstone 模。超流体中的无能隙声子就是一个 Nambu–Goldstone 模。

直观地说，如果一个对称性（连续的或分立的）被自发破缺，那么各基态必定是简并的，因为哈密顿量具有该对称性而基态没有。不同的基态由对称变换相互联系。此外，\emph{连续}对称性的破缺导致一个简并基态流形，它可以由\emph{连续}变量 $\theta_i$ 参数化。这些简并基态之间的涨落对应于无能隙激发（见图 3.7）。我们可以用 $\theta_i(\pmb{x},t)$ 场写出一个低能有效作用量，来描述无能隙模式的动力学。这个低能有效理论称为非线性 $\sigma$ 模型（nonlinear $\sigma$-model）。XY 模型是最简单的非线性 $\sigma$ 模型。后面我们会看到非线性 $\sigma$ 模型的更多例子。

\subsection{理解有限系统中的自发对称性破缺}

\begin{keypoints}
\item 量子涨落可以恢复对称性。有限系统的真实基态永远不可能破坏 $U(1)$ 对称性。然而，对于大系统，可以用许多近简并的基态构造出一个近似的对称性破缺基态。
\item 对于无限系统，由序参量 $\varphi_0\mathrm{e}^{\mathrm{i}\theta}$ 表征的态自成一个“宇宙”。它要涨落到另一个由不同序参量表征的态，需要无穷长的时间。
\end{keypoints}

% ------------------------------------------------------------
% 块边界备注：
%   起点：书页 78（ch3_a）末行止于 "The energy density (or, more precisely, the"，
%   本块正文以 %JOIN% 续书页 79 顶部半句 "density of the thermal potential) of
%   the classical ground state is"（译作"热势的密度）为"），ch3_a 应将
%   "经典基态的能量密度（或者更精确地说，"写入其 %--A-TAIL-- 区。
%   终点：书页 89（p104）以 3.3.7 节小标题下的要点框自然收尾；书页 90（p105）
%   顶部为新段落 "In Section 3.3.2, we mentioned ..."（式 (3.3.19) 起），
%   无跨页半句，故本块不需要 %--B-TAIL--，ch3_c 直接从书页 90 的新段开始。
%   本块内无 \section 级标题：3.3 节及 3.3.1、3.3.2 小节均在 ch3_a 范围
%   （书页 77–78）；本块起于 3.3.2 节中段，含 \subsection 3.3.3–3.3.7 与
%   问题 3.3.1–3.3.3。
% ------------------------------------------------------------
% 新增术语（ch3_b，待并入术语表"新增术语"区）：
%   superfluid transition 超流相变；
%   first-order / second-order phase transition 一级 / 二级相变；
%   non-analytic point 非解析点；kink 拐折；
%   degenerate ground states 简并基态；manifold 流形；
%   Gibbs energy 吉布斯能；water–vapor transition 水–水蒸气相变；
%   critical point 临界点；critical temperature 临界温度；
%   integrating out (a field) 积掉（某个场）；
%   total derivative / total time derivative 全导数 / 全时间导数；
%   wave equation 波动方程；dispersion 色散；
%   compressibility 压缩率；conductivity 电导率；
%   density–current fluctuations 密度–流涨落；conservation law 守恒定律；
%   particle–wave duality 波粒二象性；quasiparticle 准粒子（已入表）；
%   phonon 声子；roton 旋子；phonon charge 声子电荷；
%   vortex ring 涡环；vortex line 涡线；source / drain 源 / 汇；
%   multipole expansion 多极展开；dipolar interaction 偶极相互作用；
%   canonical momentum 正则动量；canonical commutation relation 正则对易关系；
%   lowering / raising operator 下降 / 上升算符；
%   decoupled harmonic oscillators 退耦谐振子；collective motion 集体运动；
%   toy universe 玩具宇宙；elementary particle 基本粒子；
%   relativistic particle 相对论性粒子；
%   universal properties 普适性质；explicit symmetry breaking 显式（对称性）破缺；
%   Nambu–Goldstone mode Nambu–Goldstone 模
% ------------------------------------------------------------
